\section{Laurent power series}

\subsection{Motivation}

Fix a commutative ring $K$.

\begin{definition}
\label{def.fps.laure.binary-rep}
\lean{AlgebraicCombinatorics.FPS.Laurent.BinaryRepresentation}
\leanhelper
\leanok
A \emph{binary representation} of an integer $n$ means an essentially finite
sequence $\left(  b_{i}\right)  _{i\in\mathbb{N}}=\left(  b_{0},b_{1}%
,b_{2},\ldots\right)  \in\left\{  0,1\right\}  ^{\mathbb{N}}$ such that%
\[
n=\sum_{i\in\mathbb{N}}b_{i}2^{i}.
\]
(Recall that ``essentially finite'' means
``all but finitely many $i\in\mathbb{N}$ satisfy $b_{i}%
=0$''.)
\end{definition}

\begin{theorem}
\label{thm.fps.laure.binary-rep-uniq}
\lean{AlgebraicCombinatorics.FPS.Laurent.binaryRepresentation_exists_unique, AlgebraicCombinatorics.binaryRepresentation_unique, AlgebraicCombinatorics.binaryRepresentation_exists}
\leantarget
\leanok
Each $n\in\mathbb{N}$ has a unique binary representation.
\end{theorem}

\begin{proof}
\leanok
In the first formalization, the proof uses
the canonical bit-index representation from Mathlib.
In the second formalization, the proof proceeds
by strong induction on $n$: the least significant bit is determined
by $n \bmod 2$, and uniqueness of the remaining bits follows from
the induction hypothesis applied to $n / 2$.
\end{proof}

\begin{definition}
\label{def.fps.laure.balanced-tern}
\lean{AlgebraicCombinatorics.BalancedTernaryRepresentation, AlgebraicCombinatorics.FPS.Laurent.BalancedTernaryRep}
\leanhelper
\leanok
A \emph{balanced ternary representation} of an integer $n$ means an
essentially finite sequence $\left(  b_{i}\right)  _{i\in\mathbb{N}}=\left(
b_{0},b_{1},b_{2},\ldots\right)  \in\left\{  0,1,-1\right\}  ^{\mathbb{N}}$
such that%
\[
n=\sum_{i\in\mathbb{N}}b_{i}3^{i}.
\]
\end{definition}

\begin{theorem}
\label{thm.fps.laure.balanced-tern-rep-uniq}
\lean{AlgebraicCombinatorics.FPS.Laurent.balancedTernary_exists_unique, AlgebraicCombinatorics.balancedTernaryRepresentation_unique, AlgebraicCombinatorics.balancedTernaryRepresentation_exists}
\leantarget
\leanok
Each integer $n$ has a unique balanced ternary representation.
\end{theorem}

\begin{proof}
\leanok
\textbf{Existence.}
In the first formalization, existence is proved
by finding $k$ large enough that $|n| \le M_k = 3^0 + 3^1 + \cdots + 3^k$,
then using a cardinality/pigeonhole argument: the map from $k$-bounded
representations (of which there are $3^{k+1}$) to the interval
$\{m \in \mathbb{Z} : |m| \le M_k\}$ (which also has size $3^{k+1}$) is
injective, hence surjective.
In the second formalization, existence is proved
by strong induction on $|n|$: the least significant digit is determined
by $n \bmod 3$ (choosing from $\{-1,0,1\}$), and the remaining digits
represent $(n - d) / 3$.

\textbf{Uniqueness.}
Both formalizations prove uniqueness by showing that if the difference
$c_i = b_i - b_i'$ of two representations satisfies
$\sum c_i \cdot 3^i = 0$ with $c_i \in \{-2,\ldots,2\}$ and all but
finitely many $c_i = 0$, then all $c_i = 0$. This is proved by
induction on the highest nonzero index: $c_0 \equiv 0 \pmod{3}$ forces
$c_0 = 0$, then dividing by $3$ reduces to the induction hypothesis.
\end{proof}

\begin{lemma}[Key uniqueness lemma for base-3 digit sums]
\label{lem.fps.laure.sum-powers-three-zero}
\lean{AlgebraicCombinatorics.FPS.Laurent.btDigits_sum_unique, AlgebraicCombinatorics.sum_powers_of_three_eq_zero}
\leanhelper
\leanok
Let $c : \mathbb{N} \to \mathbb{Z}$ be an essentially finite sequence
with $c_i \in \{-2, \ldots, 2\}$ for all $i$.
If $\sum_{i} c_i \cdot 3^i = 0$, then $c_i = 0$ for all $i$.
\end{lemma}

\begin{proof}
\leanok
By induction on the highest nonzero index $k$.
Since $c_0 \equiv 0 \pmod{3}$ and $|c_0| \le 2$, we get $c_0 = 0$.
Dividing the remaining sum by $3$ gives a shifted sequence satisfying the
same hypotheses with a smaller bound, and the induction hypothesis applies.
\end{proof}

\begin{lemma}[$k$-bounded balanced ternary uniqueness]
\label{lem.fps.laure.kbounded-unique}
\lean{AlgebraicCombinatorics.FPS.Laurent.kBounded_unique, AlgebraicCombinatorics.balancedTernaryRepresentation_bounded_unique}
\leanhelper
\leanok
Each integer $n$ with $|n| \le 3^0 + 3^1 + \cdots + 3^k$ has a unique
$k$-bounded balanced ternary representation.
\end{lemma}

\begin{proof}
\leanok
In the first formalization, this uses a cardinality argument: the number
of $k$-bounded representations is $3^{k+1}$, the target interval
$\{m : |m| \le M_k\}$ also has size $3^{k+1}$, and the evaluation map
is injective (by Lemma~\ref{lem.fps.laure.sum-powers-three-zero}), hence
bijective.
In the second formalization, uniqueness follows directly from
Lemma~\ref{lem.fps.laure.sum-powers-three-zero} applied to the difference
of two representations.
\end{proof}

\begin{lemma}[Balanced ternary digit type]
\label{lem.fps.laure.bt-digit-type}
\lean{AlgebraicCombinatorics.BalancedTernaryDigit, AlgebraicCombinatorics.BalancedTernaryDigit.toInt, AlgebraicCombinatorics.BalancedTernaryDigit.toInt_injective}
\leanhelper
\leanok
A balanced ternary digit is one of $-1$, $0$, or $1$.
There is an injection from the set of balanced ternary digits
to $\mathbb{Z}$ sending each digit to its integer value.
This injection is injective, and two digits are equal if and only if
their integer values are congruent modulo $3$.
\end{lemma}

\begin{proof}
Exhaustive case analysis on the three values of a balanced ternary digit.
\end{proof}

\begin{lemma}[Existence of balanced ternary representations by strong induction]
\label{lem.fps.laure.bt-exists-inductive}
\lean{AlgebraicCombinatorics.balancedTernaryRepresentation_exists}
\leanhelper
\leanok
Every integer $n$ has a balanced ternary representation.
The proof proceeds by strong induction on $|n|$:
the least significant digit $d$ is chosen by $n \bmod 3$
(from $\{-1, 0, 1\}$), and the remaining digits represent
$(n - d)/3$, which satisfies $|(n-d)/3| < |n|$ for $n \ne 0$.
\end{lemma}

\begin{proof}
Base case: $n = 0$ uses the zero representation (all digits zero).
Inductive step: pick the digit $d$ such that $n \equiv d \pmod{3}$
with $d \in \{-1, 0, 1\}$, form $m = (n - d)/3$, and prepend $d$
to the representation of $m$ (obtained by the induction hypothesis
since $|m| < |n|$).
\end{proof}

\begin{lemma}[Partial product formula for balanced ternary]
\label{lem.fps.laure.partial-product}
\lean{AlgebraicCombinatorics.balancedTernaryPartialProduct, AlgebraicCombinatorics.geom_sum_three}
\leanhelper
\leanok
Define the partial product
$P_k = \prod_{i=0}^{k} (1 + T(3^i) + T(-3^i))$ in $K[\xpm]$.
The geometric sum identity
$2 \sum_{i=0}^{k} 3^i = 3^{k+1} - 1$
is used to establish that $P_k$ enumerates all integers
in $[-M_k, M_k]$ where $M_k = (3^{k+1}-1)/2$.
\end{lemma}

\begin{proof}
The geometric sum identity is proved by induction on $k$.
\end{proof}

\begin{lemma}[Cardinality identity: $2 M_k + 1 = 3^{k+1}$]
\label{lem.fps.laure.maxBT-eq}
\lean{AlgebraicCombinatorics.FPS.Laurent.maxBT_eq}
\leanhelper
\leanok
For $M_k = \sum_{i=0}^{k} 3^i$, we have $2 M_k + 1 = 3^{k+1}$.
\end{lemma}

\begin{proof}
\leanok
By induction on $k$.
\end{proof}

\begin{lemma}[$2 \sum_{i<k} 3^i < 3^k$]
\label{lem.fps.laure.two-sum-lt}
\lean{AlgebraicCombinatorics.FPS.Laurent.two_sum_lt_pow_three}
\leanhelper
\leanok
For all $k \ge 0$, we have $2 \sum_{i=0}^{k-1} 3^i < 3^k$.
\end{lemma}

\begin{proof}
\leanok
By induction on $k$.
\end{proof}

\begin{lemma}[{Bounded representation values lie in $[-M_k, M_k]$}]
\label{lem.fps.laure.kBoundedBTValue-mem}
\lean{AlgebraicCombinatorics.FPS.Laurent.kBoundedBTValue_mem_Icc}
\leanhelper
\leanok
The value $\sum_{i=0}^{k} f(i) \cdot 3^i$ of any $k$-bounded
balanced ternary representation $f$ (with $f(i) \in \{-1, 0, 1\}$)
lies in the interval $[-M_k, M_k]$.
\end{lemma}

\begin{proof}
\leanok
Each term $f(i) \cdot 3^i$ satisfies $-3^i \le f(i) \cdot 3^i \le 3^i$,
so the total sum is bounded by $\pm \sum_{i=0}^{k} 3^i = \pm M_k$.
\end{proof}

\begin{lemma}[Injectivity of $k$-bounded evaluation]
\label{lem.fps.laure.balanced-ternary-injective}
If two $k$-bounded balanced ternary representations have the same value,
they must be equal. The proof uses a mod-$3$ argument:
the lowest digit is determined by the value modulo $3$, and dividing
by $3$ reduces to the induction hypothesis.
\end{lemma}

\begin{proof}
By strong induction on $k$. Base case ($k = 0$): the value equals
$f(0)$, so equality of values gives equality of digits.
Inductive step: the value modulo $3$ determines $f(0)$ (since
$f(0) \in \{-1, 0, 1\}$), then dividing by $3$ gives a
$(k-1)$-bounded problem.
\end{proof}

\begin{lemma}[Conversion between bounded and finitely supported representations]
\label{lem.fps.laure.kBoundedToFinsupp}
\lean{AlgebraicCombinatorics.FPS.Laurent.kBoundedToFinsupp, AlgebraicCombinatorics.FPS.Laurent.kBoundedToFinsupp_sum}
\leanhelper
\leanok
A $k$-bounded representation $f : \{0, 1, \ldots, k\} \to \mathbb{Z}$
can be converted to a finitely supported function
$\mathbb{N} \to \mathbb{Z}$ by extending with zeros.
The weighted sum of the resulting function equals the $k$-bounded value.
\end{lemma}

\begin{proof}
\leanok
The finitely supported function is constructed by extending with zeros.
The sum identity follows by reindexing.
\end{proof}

\subsection{The space $K\left[\left[\xpm\right]\right]$}

\begin{definition}
\label{def.fps.laure.double}
\lean{AlgebraicCombinatorics.DoublyInfinitePowerSeries}
\leantarget
\leanok
Let $K\left[  \left[  \xpm\right]  \right]  $
be the $K$-module $K^{\mathbb{Z}}$ of all families $\left(  a_{n}\right)
_{n\in\mathbb{Z}}=\left(  \ldots,a_{-2},a_{-1},a_{0},a_{1},a_{2}%
,\ldots\right)  $ of elements of $K$. Its addition and its scaling are defined
entrywise:%
\begin{align*}
\left(  a_{n}\right)  _{n\in\mathbb{Z}}+\left(  b_{n}\right)  _{n\in
\mathbb{Z}}  &  =\left(  a_{n}+b_{n}\right)  _{n\in\mathbb{Z}};\\
\lambda\left(  a_{n}\right)  _{n\in\mathbb{Z}}  &  =\left(  \lambda
a_{n}\right)  _{n\in\mathbb{Z}}\ \ \ \ \ \ \ \ \ \ \text{for each }\lambda\in
K.
\end{align*}
An element of $K\left[  \left[  \xpm\right]  \right]  $ will be called a
\emph{doubly infinite power series}. We use the notation $\sum_{n\in\mathbb{Z}}a_{n}x^{n}$ for a family
$\left(  a_{n}\right)  _{n\in\mathbb{Z}}\in K\left[  \left[  \xpm\right]
\right]  $.
\end{definition}

\begin{lemma}[Single monomials in $K\llbracket \xpm \rrbracket$]
\label{lem.fps.laure.double-single}
\lean{AlgebraicCombinatorics.DoublyInfinitePowerSeries.single, AlgebraicCombinatorics.DoublyInfinitePowerSeries.coeff_single, AlgebraicCombinatorics.DoublyInfinitePowerSeries.smul_single}
\leanhelper
\leanok
For $n \in \mathbb{Z}$, the monomial $x^n \in K\llbracket \xpm \rrbracket$
is the family $(\delta_{m,n})_{m \in \mathbb{Z}}$.
Its coefficient at $m$ is $\delta_{m,n}$.
Scalar multiplication satisfies $\lambda \cdot x^n = (\lambda \delta_{m,n})_{m \in \mathbb{Z}}$.
\end{lemma}

\begin{proof}
\leanok
Immediate from the definitions.
\end{proof}

\begin{lemma}[Coefficient decomposition of doubly infinite power series]
\label{lem.fps.laure.double-decomp}
\lean{AlgebraicCombinatorics.DoublyInfinitePowerSeries.eq_coeff_at_position, AlgebraicCombinatorics.DoublyInfinitePowerSeries.coeff_eq_single_contrib}
\leanhelper
\leanok
For any $f = (a_n) \in K\llbracket \xpm \rrbracket$ and $m \in \mathbb{Z}$,
$f$ evaluated at position $m$ gives $a_m$.
Equivalently, the coefficient of $f$ at $m$ equals the scalar
by which $f$ scales the $m$-th basis vector $(\delta_{n,m})_{n \in \mathbb{Z}}$.
\end{lemma}

\begin{proof}
\leanok
Immediate from the definitions.
\end{proof}

\subsection{Laurent polynomials}

\begin{definition}
\label{def.fps.laure.laupol}
\lean{AlgebraicCombinatorics.FPS.Laurent.LaurentPoly, AlgebraicCombinatorics.FPS.Laurent.laurentPoly_mul_coeff, AlgebraicCombinatorics.FPS.Laurent.laurentPoly_X, AlgebraicCombinatorics.laurentPolynomial_one_eq_T_zero}
\leantarget
\leanok
Let $K\left[  \xpm\right]  $ be the
$K$-submodule of $K\left[  \left[  \xpm\right]  \right]  $ consisting of
all \textbf{essentially finite} families $\left(  a_{n}\right)  _{n\in
\mathbb{Z}}$. The elements of $K\left[  \xpm\right]  $ are called \emph{Laurent
polynomials} in the indeterminate $x$ over $K$.

We define a multiplication on $K\left[  \xpm\right]  $ by setting%
\[
\left(  a_{n}\right)  _{n\in\mathbb{Z}}\cdot\left(  b_{n}\right)
_{n\in\mathbb{Z}}=\left(  c_{n}\right)  _{n\in\mathbb{Z}}%
,\ \ \ \ \ \ \ \ \ \ \text{where}\ \ \ \ \ \ \ \ \ \ c_{n}=\sum_{i\in
\mathbb{Z}}a_{i}b_{n-i}.
\]
The sum $\sum_{i\in\mathbb{Z}}a_{i}b_{n-i}$ is well-defined
because it is essentially finite.

We define an element $x\in K\left[  \xpm\right]  $ by%
\[
x=\left(  \delta_{i,1}\right)  _{i\in\mathbb{Z}}.
\]

In Mathlib, Laurent polynomials are represented as
finitely supported functions $\mathbb{Z} \to K$ (the group algebra
of $\mathbb{Z}$ over $K$).
\end{definition}

\begin{theorem}
\label{thm.fps.laure.laupol-ring}
\lean{AlgebraicCombinatorics.FPS.Laurent.laurentPolynomial_commRing, AlgebraicCombinatorics.FPS.Laurent.laurentPolynomial_algebra, AlgebraicCombinatorics.FPS.Laurent.T_isUnit, AlgebraicCombinatorics.laurentPolynomial_commSemiring, AlgebraicCombinatorics.laurentPolynomial_algebra, AlgebraicCombinatorics.laurentPolynomial_T_isUnit}
\leantarget
\leanok
The $K$-module $K\left[  \xpm\right]  $,
equipped with the multiplication defined above, is a commutative
$K$-algebra. Its unity is $\left(  \delta_{i,0}\right)  _{i\in\mathbb{Z}}$.
The element $x$ is invertible in this $K$-algebra.
\end{theorem}

\begin{proof}
\leanok
The commutative ring and algebra structures are provided by Mathlib's
group algebra structure.
Invertibility of $x = T(1)$ is witnessed by $T(-1)$: we have
$T(1) \cdot T(-1) = T(1 + (-1)) = T(0) = 1$ and similarly
$T(-1) \cdot T(1) = 1$.
\end{proof}

\begin{lemma}[Unit structure of $T(n)$]
\label{lem.fps.laure.T-inv}
\lean{AlgebraicCombinatorics.laurentPolynomial_T_unit, AlgebraicCombinatorics.laurentPolynomial_T_isUnit', AlgebraicCombinatorics.laurentPolynomial_T_mul_T_neg, AlgebraicCombinatorics.laurentPolynomial_T_neg_mul_T}
\leanhelper
\leanok
For every $n \in \mathbb{Z}$, $T(n)$ is a unit in $K[\xpm]$ with
inverse $T(-n)$: we have $T(n) \cdot T(-n) = 1$ and $T(-n) \cdot T(n) = 1$.
\end{lemma}

\begin{proof}
\leanok
By $T(a) \cdot T(b) = T(a+b)$ and $T(0) = 1$.
\end{proof}

\begin{lemma}[Laurent polynomial decomposition]
\label{lem.fps.laure.laupol-decomp}
\lean{AlgebraicCombinatorics.laurentPolynomial_eq_sum, AlgebraicCombinatorics.laurentPolynomial_eq_groupAlgebra}
\leanhelper
\leanok
Every Laurent polynomial $f \in K[\xpm]$ can be written as
$f = \sum_{n \in \mathrm{supp}(f)} f(n) \cdot T(n)$.
Moreover, $K[\xpm]$ is isomorphic to the group algebra
$K[\mathbb{Z}]$.
\end{lemma}

\begin{proof}
\leanok
The decomposition follows from the fact that any finitely supported
function equals the sum of its single-term components.
The group algebra isomorphism is the identity on the
underlying finitely supported functions $\mathbb{Z} \to K$.
\end{proof}

\begin{lemma}[Multiplication by $T$ shifts coefficients]
\label{lem.fps.laure.xa}
\lean{AlgebraicCombinatorics.FPS.Laurent.T_mul_coeff, AlgebraicCombinatorics.FPS.Laurent.T_neg_one_mul_coeff}
\leantarget
\leanok
If $a = (a_n)_{n \in \mathbb{Z}}$ is a Laurent polynomial, then
$T(1) \cdot a = (a_{n-1})_{n \in \mathbb{Z}}$ and
$T(-1) \cdot a = (a_{n+1})_{n \in \mathbb{Z}}$.
\end{lemma}

\begin{proof}
\leanok
Follows from the definition of multiplication: the convolution
of a single-term function with $a$ shifts the coefficients.
\end{proof}

\begin{lemma}[Powers of $T$]
\label{lem.fps.laure.xk}
\lean{AlgebraicCombinatorics.FPS.Laurent.T_coeff_eq, AlgebraicCombinatorics.FPS.Laurent.T_one_zpow}
\leanhelper
\leanok
For each $k \in \mathbb{Z}$, the Laurent polynomial $T(k)$ has coefficient
$1$ at position $k$ and $0$ elsewhere:
$T(k)_i = \delta_{i,k}$.
Moreover, $T(1)^k = T(k)$ for all $k \in \mathbb{Z}$ (using integer powers
via the unit structure).
\end{lemma}

\begin{proof}
The coefficient characterization follows from the fact that $T(k)$ is
the family with $1$ at position $k$ and $0$ elsewhere.
The power identity $T(1)^k = T(k)$ is proved by induction on $k$
(for $k \geq 0$, using $T(a) \cdot T(b) = T(a+b)$; for $k < 0$,
using the inverse $T(-1)$).
\end{proof}

\begin{lemma}[Coefficient of $T(k)$ at position $n$]
\label{lem.fps.laure.xk.coeff}
\lean{AlgebraicCombinatorics.FPS.Laurent.coeff_powerSeriesToLaurent}
\leanhelper
\leanok
For any $k \in \mathbb{Z}$ and $n \in \mathbb{Z}$, the coefficient
$[x^n] T(k)$ equals $\delta_{n,k}$, i.e., $1$ if $n = k$ and $0$ otherwise.
For power series embedded into Laurent series, the coefficient at
non-negative index $n \in \mathbb{N}$ is preserved.
\end{lemma}

\begin{proof}
\leanok
Follows from the fact that $T(k)$ is the family $(\delta_{n,k})_{n \in \mathbb{Z}}$
and the definition of coefficients.
\end{proof}

\begin{lemma}[Right-multiplication by $T$ on Laurent series]
\label{lem.fps.laure.xa.right}
\lean{AlgebraicCombinatorics.FPS.Laurent.laurentSeries_eq_hsum_single}
\leanhelper
\leanok
A Laurent series $f \in K((x))$ equals the summable-family sum of the
monomials $f_n \cdot x^n$ (each being the family $(\delta_{m,n} \cdot f_n)_{m \in \mathbb{Z}}$).
This constructs an explicit summable family whose sum recovers $f$.
\end{lemma}

\begin{proof}
\leanok
Construct the summable family by sending $n \in \mathbb{Z}$ to
the monomial $f_n \cdot x^n$.
The support at any given position is a singleton (or empty), hence finite.
Checking coefficients: at position $m$, only the $m$-th monomial contributes.
\end{proof}

\begin{proposition}
\label{prop.fps.laure.a=sumaixi}
\lean{AlgebraicCombinatorics.FPS.Laurent.eq_sum_T, AlgebraicCombinatorics.FPS.Laurent.laurentSeries_eq_hsum_single, AlgebraicCombinatorics.laurentPolynomial_eq_sum}
\leantarget
\leanok
Any doubly infinite power series $a=\left(
a_{i}\right)  _{i\in\mathbb{Z}}\in K\left[  \left[  \xpm\right]  \right]  $
satisfies%
\[
a=\sum_{i\in\mathbb{Z}}a_{i}x^{i}.
\]
Here, the powers $x^{i}$ are taken in the Laurent polynomial ring $K\left[
\xpm\right]  $, but the infinite sum $\sum_{i\in\mathbb{Z}}a_{i}x^{i}$ is
taken in the $K$-module $K\left[  \left[  \xpm\right]  \right]  $.

For Laurent polynomials, the sum is finite and the identity holds in
$K[\xpm]$ itself.
For Laurent series, the identity is formalized using a summable family construction.
\end{proposition}

\begin{proof}
For Laurent polynomials, this follows from the decomposition of
a finitely supported function as a sum of its single-term components.
For Laurent series, the proof constructs a summable family of
monomials and shows that its sum equals the original
series by checking coefficients: at each position $n$, only the $n$-th
monomial contributes (using the coefficient characterization of $x^k$
from Lemma~\ref{lem.fps.laure.xk} and Lemma~\ref{lem.fps.laure.xk.coeff}).
\end{proof}

\subsection{Laurent series}

\begin{definition}
\label{def.fps.laure.lauser}
\lean{AlgebraicCombinatorics.DoublyInfinitePowerSeries.IsLaurentSeries, AlgebraicCombinatorics.FPS.Laurent.laurentSeries_commRing}
\leantarget
\leanok
We let $K\left(  \left(  x\right)  \right)  $ be
the subset of $K\left[  \left[  \xpm\right]  \right]  $ consisting of all
families $\left(  a_{i}\right)  _{i\in\mathbb{Z}}\in K\left[  \left[  x^{\pm
}\right]  \right]  $ such that the sequence $\left(  a_{-1},a_{-2}%
,a_{-3},\ldots\right)  $ is essentially finite -- i.e., such that all
sufficiently low $i\in\mathbb{Z}$ satisfy $a_{i}=0$.

The elements of $K\left(  \left(  x\right)  \right)  $ are called
\emph{Laurent series} in one indeterminate $x$ over $K$.

In Mathlib, Laurent series are implemented as Hahn series
over $\mathbb{Z}$ with coefficients in $K$,
which are functions $\mathbb{Z} \to K$ whose support is well-founded
(equivalently, bounded below). The formalization also provides
a predicate on doubly infinite power series that captures
this textbook definition.
\end{definition}

\begin{lemma}[Mathlib's Laurent series satisfy the textbook definition]
\label{lem.fps.laure.mathlib-agrees}
\lean{AlgebraicCombinatorics.DoublyInfinitePowerSeries.isLaurentSeries_ofLaurentSeries, AlgebraicCombinatorics.DoublyInfinitePowerSeries.toLaurentSeries_ofLaurentSeries, AlgebraicCombinatorics.DoublyInfinitePowerSeries.ofLaurentSeries_toLaurentSeries}
\leanhelper
\leanok
Mathlib's Laurent series satisfy the textbook definition
(the sequence of negative-indexed coefficients is essentially finite).
Moreover, the conversions between doubly infinite power series
satisfying this condition and Mathlib's Laurent series
are inverse to each other.
\end{lemma}

\begin{proof}
\leanok
The forward direction uses the well-foundedness of the support:
a Hahn series has a minimum support element (if nonempty), and all
indices below it have zero coefficient.
The round-trip identities hold by direct computation on coefficients.
\end{proof}

\begin{lemma}[Bounded-below subsets of $\mathbb{Z}$ are partially well-ordered]
\label{lem.fps.laure.isPWO-bddBelow}
\lean{AlgebraicCombinatorics.DoublyInfinitePowerSeries.Int.isPWO_of_bddBelow}
\leanhelper
\leanok
If $s \subseteq \mathbb{Z}$ is bounded below, then $s$ is partially
well-ordered.
\end{lemma}

\begin{proof}
\leanok
Given a sequence in $s$, shift it to $\mathbb{N}$ by subtracting the
lower bound. By Ramsey's theorem (monotone subsequence in $\mathbb{N}$),
extract a monotone subsequence. A strictly decreasing sequence in
$\mathbb{N}$ is impossible, so the monotone case gives the result.
\end{proof}

\begin{theorem}
\label{thm.fps.laure.lauser-ring}
\lean{AlgebraicCombinatorics.FPS.Laurent.laurentSeries_commRing, AlgebraicCombinatorics.FPS.Laurent.laurentSeries_algebra, AlgebraicCombinatorics.FPS.Laurent.LaurentSeries_X_isUnit}
\leantarget
\leanok
The subset $K\left(  \left(  x\right)  \right)  $
is a $K$-submodule of $K\left[  \left[  \xpm\right]  \right]  $.
It has a multiplication given by the same convolution rule as
$K\left[  \xpm\right]  $: namely,%
\[
\left(  a_{n}\right)  _{n\in\mathbb{Z}}\cdot\left(  b_{n}\right)
_{n\in\mathbb{Z}}=\left(  c_{n}\right)  _{n\in\mathbb{Z}}%
,\ \ \ \ \ \ \ \ \ \ \text{where}\ \ \ \ \ \ \ \ \ \ c_{n}=\sum_{i\in
\mathbb{Z}}a_{i}b_{n-i}.
\]
The sum $\sum_{i\in\mathbb{Z}}a_{i}b_{n-i}$ is well-defined because
for sufficiently low $i$ we have $a_i = 0$, and for sufficiently high $i$
we have $b_{n-i} = 0$.

Equipped with this multiplication, $K\left(
\left(  x\right)  \right)  $ is a commutative $K$-algebra with unity
$\left(  \delta_{i,0}\right)  _{i\in\mathbb{Z}}$.
The element $x$ is invertible in this algebra.
\end{theorem}

\begin{proof}
The commutative ring and algebra structures follow from Mathlib's
Laurent series ring structure, using the agreement between our definition
and Mathlib's (Lemma~\ref{lem.fps.laure.mathlib-agrees}).
Well-definedness of the product
follows from the partial well-orderedness of the support
and the order multiplicativity (Lemma~\ref{lem.fps.laure.order-mul}).

The algebra structure extends the FPS algebra via
the power series embedding (Lemma~\ref{lem.fps.laure.ps-embed}) and
includes Laurent polynomials via the Laurent polynomial embedding
(Lemma~\ref{lem.fps.laure.laupol-embed}).

Invertibility of $x$: the element $x = x^1$ has inverse
$x^{-1}$, since $x^1 \cdot x^{-1} = x^{1+(-1)} = x^0 = 1$.
\end{proof}

\begin{lemma}[Order of product]
\label{lem.fps.laure.order-mul}
\lean{AlgebraicCombinatorics.FPS.Laurent.order_mul_ge, AlgebraicCombinatorics.FPS.Laurent.coeff_mul_eq_zero_of_lt_order}
\leanhelper
\leanok
If $f$ and $g$ are Laurent series, then
$\operatorname{ord}(f) + \operatorname{ord}(g) \le \operatorname{ord}(f \cdot g)$.
In particular, the coefficients of $f \cdot g$ vanish below
$\operatorname{ord}(f) + \operatorname{ord}(g)$.
\end{lemma}

\begin{proof}
\leanok
This is a standard property of Hahn series from Mathlib.
\end{proof}

\subsection{Embeddings}

\begin{lemma}[Power series embed into Laurent series]
\label{lem.fps.laure.ps-embed}
\lean{AlgebraicCombinatorics.FPS.Laurent.powerSeriesToLaurent, AlgebraicCombinatorics.FPS.Laurent.powerSeriesToLaurent_injective}
\leanhelper
\leanok
There is an injective ring homomorphism
$K\left[\left[x\right]\right] \hookrightarrow K\left(\left(x\right)\right)$
that preserves coefficients.
\end{lemma}

\begin{proof}
\leanok
This is provided by Mathlib's embedding of power series into
Hahn series, which is injective.
\end{proof}

\begin{lemma}[Laurent polynomials embed into Laurent series]
\label{lem.fps.laure.laupol-embed}
\lean{AlgebraicCombinatorics.FPS.Laurent.laurentPolynomialToSeries, AlgebraicCombinatorics.FPS.Laurent.laurentPolynomialToSeries_mul, AlgebraicCombinatorics.laurentPolyToSeries, AlgebraicCombinatorics.laurentPolyToSeries_injective}
\leanhelper
\leanok
There is an injective ring homomorphism
$K\left[\xpm\right] \hookrightarrow K\left(\left(x\right)\right)$
that preserves coefficients.
\end{lemma}

\begin{proof}
\leanok
In the first formalization, the map sends a finitely supported function
$p : \mathbb{Z} \to K$ to the Laurent series with the same coefficient function
(whose support, being finite, is automatically bounded below).
Additivity and multiplicativity are checked by comparing coefficients.

In the second formalization, the map is constructed as a ring homomorphism
using the universal property of the group algebra, sending
$n \in \mathbb{Z}$ to the monomial $x^n$ in $K((x))$.
\end{proof}

\begin{lemma}[Embedding preserves addition]
\label{lem.fps.laure.laupol-embed-add}
\lean{AlgebraicCombinatorics.FPS.Laurent.laurentPolynomialToSeries_add}
\leanhelper
\leanok
The embedding $K[\xpm] \hookrightarrow K((x))$ preserves addition:
for Laurent polynomials $p, q$, the image of $p + q$ equals
the sum of their images.
\end{lemma}

\begin{proof}
\leanok
By checking coefficients: both sides agree at every position $n \in \mathbb{Z}$.
\end{proof}

\begin{lemma}[Embedding preserves $0$ and $1$]
\label{lem.fps.laure.laupol-embed-one}
\lean{AlgebraicCombinatorics.FPS.Laurent.laurentPolynomialToSeries_zero, AlgebraicCombinatorics.FPS.Laurent.laurentPolynomialToSeries_one}
\leanhelper
\leanok
The embedding sends $0 \in K[\xpm]$ to $0 \in K((x))$ and
$1 \in K[\xpm]$ to $1 \in K((x))$.
\end{lemma}

\begin{proof}
\leanok
By checking coefficients at each position.
\end{proof}

\begin{lemma}[Embedding preserves coefficients]
\label{lem.fps.laure.laupol-embed-coeff}
\lean{AlgebraicCombinatorics.FPS.Laurent.laurentPolynomialToSeries_coeff, AlgebraicCombinatorics.laurentPolyToSeries_coeff, AlgebraicCombinatorics.laurentPolyToSeries_single}
\leanhelper
\leanok
The embedding preserves coefficients: for any Laurent polynomial $p$
and $n \in \mathbb{Z}$, the $n$-th coefficient of the image equals
$p(n)$.
Moreover, $T(n)$ maps to the monomial $x^n$ in $K((x))$.
\end{lemma}

\begin{proof}
\leanok
Immediate from the definition of the embedding.
\end{proof}

\subsection{The partial product formula}

\begin{lemma}[Geometric identity for balanced ternary factors]
\label{lem.fps.laure.one-plus-T}
\lean{AlgebraicCombinatorics.one_plus_T_plus_Tinv_eq}
\leanhelper
\leanok
In the Laurent polynomial ring,
\[
(1 + T(1) + T(-1)) \cdot (T(1) \cdot (1 - T(1))) = 1 - T(3).
\]
\end{lemma}

\begin{proof}
\leanok
By direct algebraic computation using $T(a) \cdot T(b) = T(a+b)$
and the ring axioms.
\end{proof}

\begin{lemma}[Invertibility of $1 - x^i$ for $i > 0$]
\label{lem.fps.laure.one-sub-xpow-unit}
\lean{AlgebraicCombinatorics.one_sub_X_pow_isUnit}
\leanhelper
\leanok
For any positive integer $i$, the power series $1 - x^i$ is invertible
in $K\left[\left[x\right]\right]$ and hence also in
$K\left(\left(x\right)\right)$.
\end{lemma}

\begin{proof}
\leanok
The constant coefficient of $1 - x^i$ is $1$, which is a unit.
By the invertibility criterion for power series (a power series is
invertible if and only if its constant coefficient is a unit),
$1 - x^i$ is invertible.
\end{proof}

\subsection{Laurent polynomial action and torsion}

\begin{lemma}[Laurent polynomial action on doubly infinite power series]
\label{lem.fps.laure.laupol-smul}
\lean{AlgebraicCombinatorics.laurentPolynomialSmul}
\leanhelper
\leanok
A Laurent polynomial $p \in K[\xpm]$ acts on a doubly infinite
power series $f \in K\llbracket \xpm \rrbracket$ by the
convolution formula:
\[
(p \cdot f)_n = \sum_{m \in \mathrm{supp}(p)} p(m) \cdot f(n - m).
\]
\end{lemma}

\begin{proof}
Well-definedness follows from the finite support of $p$, which
makes the sum finite.
\end{proof}

\begin{lemma}[Torsion example]
\label{lem.fps.laure.torsion-example}
\lean{AlgebraicCombinatorics.torsion_example}
\leanhelper
\leanok
The constant series $\mathbf{1} = (1, 1, 1, \ldots) \in K\llbracket \xpm \rrbracket$
satisfies $(1 - x) \cdot \mathbf{1} = 0$.
This shows that the $K[\xpm]$-module $K\llbracket \xpm \rrbracket$
has torsion: a nonzero module element is annihilated by a nonzero
ring element.
\end{lemma}

\begin{proof}
The coefficient at position $n$ of $(1 - x) \cdot \mathbf{1}$ is
$1 \cdot 1 + (-1) \cdot 1 = 0$ (the contribution from position $n$
minus the contribution from position $n - 1$, both equal to $1$).
\end{proof}

\subsection{Equivalence of balanced ternary formulations}

\begin{lemma}[Equivalence between the two balanced ternary representations]
\label{lem.fps.laure.bt-equiv}
\lean{AlgebraicCombinatorics.FPS.Laurent.balancedTernary_equiv, AlgebraicCombinatorics.FPS.Laurent.BTRep.toFinsupp, AlgebraicCombinatorics.FPS.Laurent.BalancedTernaryRep.toInductive}
\leanhelper
\leanok
The two formulations of balanced ternary representations
(the finitely-supported-function formulation and
the digit-sequence formulation) are equivalent:
for each $n \in \mathbb{Z}$, the two types of balanced ternary
representations of $n$ are in bijection.
The conversions are mutual inverses (up to the digit representation).
\end{lemma}

\begin{proof}
The forward direction converts a finitely supported function with values in
$\{-1, 0, 1\}$ to a digit sequence $\mathbb{N} \to \{-1, 0, 1\}$
via the natural identification.
The backward direction converts a digit sequence
to a finitely supported function via the integer embedding.
Round-trip identities follow from the fact that the conversions between
digits and integers are mutually inverse.
\end{proof}

\begin{lemma}[Bounded representation predicate and sum]
\label{lem.fps.laure.bt-bounded-helpers}
\lean{AlgebraicCombinatorics.BalancedTernaryRepresentation.isBounded, AlgebraicCombinatorics.bounded_sum_eq}
\leanhelper
\leanok
A balanced ternary representation is \emph{$k$-bounded} if all digits
at positions $> k$ are zero.
For a $k$-bounded representation of $n$, the value equals the finite sum
$\sum_{i=0}^{k} d_i \cdot 3^i$.
\end{lemma}

\begin{proof}
\leanok
The predicate is $\forall i > k,\; d_i = 0$. The sum identity
follows from rewriting the infinite sum as a finite sum over
$\{0, \ldots, k\}$ using the vanishing of higher digits.
\end{proof}
